\textbf{Problème 5.13 (Dujella, 2026).} Pour un entier positif $k$, notons $n_1(k)$ le plus petit entier non nul (en valeur absolue) pour lequel il existe un triplet $\{a,b,c\}$ d'entiers non nuls et un ensemble $N$ d'entiers, tels que $|N|=k$, $n_1(k) \in N$, et $\{a,b,c\}$ soit un $D(n)$-triplet pour tout $n \in N$.
\begin{enumerate}
\item Est-il vrai que $n_1(5)=36$ ?
\item Que peut-on dire (en admettant certaines conjectures plausibles) du comportement asymptotique de $|n_1(k)|$ pour $k$ grand ?
\end{enumerate}

\medskip

\textit{Remarque.} Sadowski (2026) a répondu à la question (a) en montrant que $n_1(5) \leqslant 4$. Il a également montré que $n_1(6) \leqslant 4$, avec l'exemple suivant : $\{a,b,c\}=\{30,1056,65520\}$, $N=\{4,266436,15248601,21479364,63936324,1037903409\}$. De plus, $n_1(7) \leqslant 36$ avec $\{a,b,c\}=\{8,1620,257040\}$, $N=\{36,857529,1470564,200322756,238596849,452017161,16308809476\}$.
